\section{Conclusions}

The use of a complexity regularizer to control overfitting is both
standard and well-understood in machine learning. While the use of
such a regularizer introduces a trade-off --- goodness of fit
vs. model complexity --- it does not introduce a {\em tension\/},
because complexity regularization is always in service of improving
generalization, and is usually not a goal in its own right.

In contrast, in this work we have studied a variety of {\em
  fairness\/} regularizers for regression problems, and applied them
to data sets in which fairness is not subservient to generalization,
but is instead a first-order consideration.  Our empirical study has
demonstrated that the choice of fairness regularizer (group,
individual, hybrid, or other) and the particular data set can have
qualitative effects on the trade-off between accuracy and
fairness. Combined with recent theoretical
results~\cite{KleinbergMR16,Chouldechova17,FriedlerSV16} that also
highlight the incompatibility of various fairness measures, our
results highlight the care that must be taken by practitioners in
defining the type of fairness they care about for their particular
application, and in determining the appropriate balance between
predictive accuracy and fairness.

